\begin{proof}[Proof of \cref{obj:thm:five-output-upper-bound}]
\begin{enumerate}
\item
Use the record and pattern indexing of \cref{obj:synth_8}. For every \(S\in\mathcal S\), define
\[
\begin{aligned}
b_{S,j}
&:=1+(e^\varepsilon-1)\mathbf 1_{\{j\in S\}},
&&j\in\{0,1,2,3\},\\
\mathbf b_S&:=(b_{S,0},b_{S,1},b_{S,2},b_{S,3})^\top,\\
g_S
&:=(e^\varepsilon-1)
\begin{pmatrix}
(1-p)(\mathbf 1_{\{1\in S\}}-\mathbf 1_{\{0\in S\}})\\
p(\mathbf 1_{\{3\in S\}}-\mathbf 1_{\{2\in S\}})
\end{pmatrix},\\
h_S(\theta)&:=\sum_{j=0}^{3}\pi_{\theta,j}b_{S,j},
\qquad
v(t):=(t,t+1)^\top,\quad t\in\mathbb R.
\end{aligned}
\]
The record probabilities are
\[
(\pi_{\theta,0},\pi_{\theta,1},\pi_{\theta,2},\pi_{\theta,3})
=
\bigl((1-p)(1-\mu_0),(1-p)\mu_0,
p(1-\mu_1),p\mu_1\bigr).
\]
The interior hypotheses make each probability positive, and
\(b_{S,j}\ge1\). Consequently,
\[
h_S(\theta)>0\qquad(S\in\mathcal S).
\]
Write
\[
\begin{aligned}
\mathcal A_\varepsilon
&=\left\{\beta\in\mathbb R^{\mathcal S}:
\beta_S\ge0,\ 
\sum_{S\in\mathcal S}\beta_S\mathbf b_S=\mathbf1
\right\},
\qquad \mathbf1:=(1,1,1,1)^\top,\\
f_S(t)&:=\frac{(g_S^\top v(t))^2}{h_S(\theta)},
\qquad
F_\theta(\beta,t)=\sum_{S\in\mathcal S}\beta_Sf_S(t),\\
E(t)&:=\sup_{\beta\in\mathcal A_\varepsilon}F_\theta(\beta,t),
\qquad t\in\mathbb R.
\end{aligned}
\]
By \cref{obj:thm:finite-oracle}, there is \(t^*\in\mathbb R\) such that
\[
E(t^*)=\inf_{t\in\mathbb R}E(t)
=J^*(\theta,p,\varepsilon).
\]
% lean: CausalSmith.Stat.LdpAteEfficiencySurface.exists_sparse_staircase_saddle

\item
We first obtain a weight vector attaining the profiled optimum.
The complementary patterns \(S_1=\{0\}\) and \(S_{14}=\{1,2,3\}\)
give the feasible vector
\[
\alpha^{\mathrm{base}}_S:=
\begin{cases}
(e^\varepsilon+1)^{-1},&S\in\{S_1,S_{14}\},\\
0,&\text{otherwise},
\end{cases}
\qquad
\sum_{S\in\mathcal S}\alpha^{\mathrm{base}}_S\mathbf b_S
=
\frac{\mathbf b_{S_1}+\mathbf b_{S_{14}}}{e^\varepsilon+1}
=\mathbf1.
\]
Thus \(\mathcal A_\varepsilon\) is nonempty. It is closed, since its
constraints are finitely many continuous equalities and inequalities.
Moreover, for every \(\beta\in\mathcal A_\varepsilon\) and \(S\in\mathcal S\),
\[
0\le\beta_S
\le\sum_{T\in\mathcal S}\beta_T
\le\sum_{T\in\mathcal S}\beta_Tb_{T,0}
=1.
\]
Hence \(\mathcal A_\varepsilon\) is compact.

Define the profiled objective on this set by
\[
\Phi(\beta):=\inf_{t\in\mathbb R}F_\theta(\beta,t),
\qquad \beta\in\mathcal A_\varepsilon.
\]
Nonnegativity of the weights and positivity of the denominators give
\[
0\le\Phi(\beta)\le F_\theta(\beta,0)<\infty.
\]
For each fixed \(t\), \(F_\theta(\,\cdot\,,t)\) is continuous.
Its infimum over \(t\) is therefore upper semicontinuous.
The maximum of \(\Phi\) on the nonempty compact feasible set is attained.
Using the optimized value in \cref{obj:def:staircase-saddle}, choose
\(\alpha^{\mathrm{opt}}\in\mathcal A_\varepsilon\) with
\[
\Phi(\alpha^{\mathrm{opt}})
=J^*(\theta,p,\varepsilon).
\]
In particular,
\[
J^*(\theta,p,\varepsilon)
\le F_\theta(\alpha^{\mathrm{opt}},t)
\qquad(t\in\mathbb R).
\]
% lean: CausalSmith.Stat.LdpAteEfficiencySurface.Jstar_exists_optimal_weight

\item
Continuity and compactness also supply a maximizer at the fixed
coordinate \(t^*\): choose
\[
\alpha^{\mathrm{env}}
\in\operatorname*{arg\,max}_{\beta\in\mathcal A_\varepsilon}
F_\theta(\beta,t^*).
\]
Then
\[
E(t^*)=F_\theta(\alpha^{\mathrm{env}},t^*),
\]
and the preceding lower bound yields
\[
J^*(\theta,p,\varepsilon)
\le F_\theta(\alpha^{\mathrm{opt}},t^*)
\le F_\theta(\alpha^{\mathrm{env}},t^*)
=E(t^*)
=J^*(\theta,p,\varepsilon).
\]
Therefore
\[
F_\theta(\alpha^{\mathrm{opt}},t^*)
=J^*(\theta,p,\varepsilon),
\qquad
F_\theta(\beta,t^*)
\le F_\theta(\alpha^{\mathrm{opt}},t^*)
\quad(\beta\in\mathcal A_\varepsilon).
\]
% lean: CausalSmith.Stat.LdpAteEfficiencySurface.informationObjective_exists_maximizer

\item
Define, for \(S\in\mathcal S\), \(t\in\mathbb R\), and
\(\beta\in\mathbb R^{\mathcal S}\),
\[
\begin{aligned}
\dot f_S(t)
&:=\frac{2(g_S^\top v(t))(g_{S,0}+g_{S,1})}{h_S(\theta)},\\
D_\beta(t)&:=\sum_{S\in\mathcal S}\beta_S\dot f_S(t),\\
C_\beta&:=\sum_{S\in\mathcal S}
\beta_S\frac{(g_{S,0}+g_{S,1})^2}{h_S(\theta)}.
\end{aligned}
\]
Expanding each square gives the identity
\[
F_\theta(\beta,t+\lambda)-F_\theta(\beta,t)
=
\lambda D_\beta(t)+\lambda^2C_\beta
\qquad(t,\lambda\in\mathbb R).
\]
The preceding step shows that \(t^*\) minimizes
\(F_\theta(\alpha^{\mathrm{opt}},\,\cdot\,)\).
If \(D_{\alpha^{\mathrm{opt}}}(t^*)\ne0\), take
\[
\lambda_{\mathrm{stat}}
:=
-\frac{D_{\alpha^{\mathrm{opt}}}(t^*)}
{2(|C_{\alpha^{\mathrm{opt}}}|+1)}.
\]
The quadratic identity then gives
\[
\begin{aligned}
0
&\le
F_\theta(\alpha^{\mathrm{opt}},t^*+\lambda_{\mathrm{stat}})
-F_\theta(\alpha^{\mathrm{opt}},t^*)\\
&=
\frac{D_{\alpha^{\mathrm{opt}}}(t^*)^2}
{4(|C_{\alpha^{\mathrm{opt}}}|+1)^2}
\left[
C_{\alpha^{\mathrm{opt}}}
-2(|C_{\alpha^{\mathrm{opt}}}|+1)
\right]
<0,
\end{aligned}
\]
where the strict inequality uses
\(C_{\alpha^{\mathrm{opt}}}\le|C_{\alpha^{\mathrm{opt}}}|\).
This contradiction proves
\[
D_{\alpha^{\mathrm{opt}}}(t^*)=0.
\]
% lean: CausalSmith.Stat.LdpAteEfficiencySurface.informationObjective_stationary_of_min

\item
For every real weight vector, define its indexed support by
\[
\operatorname{supp}(\beta)
:=\{S\in\mathcal S:\beta_S\ne0\},
\qquad \beta\in\mathbb R^{\mathcal S}.
\]
Consider
\[
\mathcal W:=
\left\{\beta\in\mathcal A_\varepsilon:
D_\beta(t^*)=D_{\alpha^{\mathrm{opt}}}(t^*),\
F_\theta(\beta,t^*)=F_\theta(\alpha^{\mathrm{opt}},t^*)
\right\}.
\]
This set contains \(\alpha^{\mathrm{opt}}\).
Choose a member of minimum support cardinality:
\[
\alpha^{\mathrm{sp}}
\in\operatorname*{arg\,min}_{\beta\in\mathcal W}
|\operatorname{supp}(\beta)|,
\qquad
\mathcal K_{\mathrm{sp}}:=\operatorname{supp}(\alpha^{\mathrm{sp}}).
\]
Such a choice exists because support cardinalities are nonnegative
integers. By construction,
\[
D_{\alpha^{\mathrm{sp}}}(t^*)=0,
\qquad
F_\theta(\alpha^{\mathrm{sp}},t^*)=J^*(\theta,p,\varepsilon),
\]
and
\[
F_\theta(\beta,t^*)\le F_\theta(\alpha^{\mathrm{sp}},t^*)
\qquad(\beta\in\mathcal A_\varepsilon).
\]

Suppose that \(|\mathcal K_{\mathrm{sp}}|>5\).
The more than five columns
\[
\begin{pmatrix}\mathbf b_S\\ \dot f_S(t^*)\end{pmatrix}
\in\mathbb R^5,
\qquad S\in\mathcal K_{\mathrm{sp}},
\]
are linearly dependent. Extending their dependence coefficients by zero
outside \(\mathcal K_{\mathrm{sp}}\) gives
\[
\begin{gathered}
\Delta=(\Delta_S)_{S\in\mathcal S}
\in\mathbb R^{\mathcal S}\setminus\{0\},
\qquad
\Delta_S=0\quad(S\notin\mathcal K_{\mathrm{sp}}),\\
\sum_{S\in\mathcal S}\Delta_S\mathbf b_S=0,
\qquad
\sum_{S\in\mathcal S}\Delta_S\dot f_S(t^*)=0.
\end{gathered}
\]
% lean: Causalean.Mathlib.Optimization.exists_kernel_direction

\item
The objective moment of this direction also vanishes.
To see this, define
\[
\delta_S:=
\begin{cases}
1,&\alpha^{\mathrm{sp}}_S=0,\\
\displaystyle\frac{\alpha^{\mathrm{sp}}_S}{|\Delta_S|+1},
&\alpha^{\mathrm{sp}}_S>0,
\end{cases}
\qquad
\delta:=\min_{S\in\mathcal S}\delta_S.
\]
Each \(\delta_S\) is positive. Since \(\mathcal S\) is finite and
nonempty, and \(\Delta_S=0\) wherever \(\alpha^{\mathrm{sp}}_S=0\),
\[
\delta>0,
\qquad
\delta|\Delta_S|\le\alpha^{\mathrm{sp}}_S
\quad(S\in\mathcal S).
\]
Thus both perturbations are feasible:
\[
\alpha^{\mathrm{sp}}\pm\delta\Delta\ge0,
\qquad
\sum_{S\in\mathcal S}
(\alpha^{\mathrm{sp}}_S\pm\delta\Delta_S)\mathbf b_S
=\mathbf1.
\]
Their objective values satisfy
\[
F_\theta(\alpha^{\mathrm{sp}}\pm\delta\Delta,t^*)
=
F_\theta(\alpha^{\mathrm{sp}},t^*)
\pm\delta\sum_{S\in\mathcal S}\Delta_Sf_S(t^*)
\le F_\theta(\alpha^{\mathrm{sp}},t^*).
\]
Using both signs and \(\delta>0\) proves
\[
\sum_{S\in\mathcal S}\Delta_Sf_S(t^*)=0.
\]
% lean: Causalean.Mathlib.Optimization.objective_orthogonal_of_maximal

\item
We now move to a boundary point with smaller support.
Choose \(S_{\mathrm{pivot}}\in\mathcal S\) with
\(\Delta_{S_{\mathrm{pivot}}}\ne0\), and define
\[
\sigma:=
\begin{cases}
1,&\Delta_{S_{\mathrm{pivot}}}>0,\\
-1,&\Delta_{S_{\mathrm{pivot}}}<0,
\end{cases}
\qquad
\Delta^{\mathrm{mov}}_S:=\sigma\Delta_S,
\qquad
\mathcal K_+:=\{S\in\mathcal S:\Delta^{\mathrm{mov}}_S>0\}.
\]
Then \(\mathcal K_+\) is nonempty. Define
\[
\lambda_{\mathrm{bdry}}
:=
\min_{S\in\mathcal K_+}
\frac{\alpha^{\mathrm{sp}}_S}{\Delta^{\mathrm{mov}}_S}
\ge0,
\qquad
\alpha^{\mathrm{bdry}}_S
:=
\alpha^{\mathrm{sp}}_S
-\lambda_{\mathrm{bdry}}\Delta^{\mathrm{mov}}_S.
\]
For \(S\in\mathcal K_+\), the defining minimum gives
\[
\lambda_{\mathrm{bdry}}\Delta^{\mathrm{mov}}_S
\le\alpha^{\mathrm{sp}}_S.
\]
For \(S\notin\mathcal K_+\),
\[
\alpha^{\mathrm{bdry}}_S
\ge\alpha^{\mathrm{sp}}_S\ge0.
\]
Hence \(\alpha^{\mathrm{bdry}}\ge0\).
Choose an index attaining the minimum, so that
\[
S_{\mathrm{hit}}\in\mathcal K_+,
\qquad
\lambda_{\mathrm{bdry}}
=
\frac{\alpha^{\mathrm{sp}}_{S_{\mathrm{hit}}}}
{\Delta^{\mathrm{mov}}_{S_{\mathrm{hit}}}}.
\]
The support condition on \(\Delta\) implies
\(\alpha^{\mathrm{sp}}_{S_{\mathrm{hit}}}>0\), whereas
\[
\alpha^{\mathrm{bdry}}_{S_{\mathrm{hit}}}=0.
\]
Coordinates outside \(\mathcal K_{\mathrm{sp}}\) remain zero.
Consequently,
\[
\operatorname{supp}(\alpha^{\mathrm{bdry}})
\subsetneq\operatorname{supp}(\alpha^{\mathrm{sp}}).
\]

The three vanishing moments established above give
\[
\begin{aligned}
\sum_{S\in\mathcal S}\alpha^{\mathrm{bdry}}_S\mathbf b_S
&=\mathbf1,\\
D_{\alpha^{\mathrm{bdry}}}(t^*)
&=D_{\alpha^{\mathrm{sp}}}(t^*)=0,\\
F_\theta(\alpha^{\mathrm{bdry}},t^*)
&=F_\theta(\alpha^{\mathrm{sp}},t^*)
=J^*(\theta,p,\varepsilon).
\end{aligned}
\]
Thus \(\alpha^{\mathrm{bdry}}\in\mathcal W\), contradicting the
minimum-support choice. We conclude that
\[
|\operatorname{supp}(\alpha^{\mathrm{sp}})|\le5.
\]
% lean: Causalean.Mathlib.Optimization.exists_sparse_maximizer

\item
Nonnegative weights and positive denominators imply
\[
C_{\alpha^{\mathrm{sp}}}
=
\sum_{S\in\mathcal S}
\alpha^{\mathrm{sp}}_S
\frac{(g_{S,0}+g_{S,1})^2}{h_S(\theta)}
\ge0.
\]
Since \(D_{\alpha^{\mathrm{sp}}}(t^*)=0\), the quadratic identity yields
\[
F_\theta(\alpha^{\mathrm{sp}},t)
-F_\theta(\alpha^{\mathrm{sp}},t^*)
=
(t-t^*)^2C_{\alpha^{\mathrm{sp}}}\ge0
\qquad(t\in\mathbb R).
\]
Therefore \(t^*\) is a global minimizer, including when the quadratic
coefficient is zero, and
\[
\inf_{t\in\mathbb R}F_\theta(\alpha^{\mathrm{sp}},t)
=
F_\theta(\alpha^{\mathrm{sp}},t^*)
=
J^*(\theta,p,\varepsilon).
\]
% lean: CausalSmith.Stat.LdpAteEfficiencySurface.informationObjective_min_of_stationary

\item
Finally, set
\[
\alpha:=\alpha^{\mathrm{sp}},
\qquad
\mathcal S_{\mathrm{out}}
:=
\{S\in\mathcal S:(P_\theta Q_\alpha)(\{S\})\ne0\}.
\]
For these feasible weights, the staircase construction gives
\[
Q_\alpha(\{S\}\mid x_j)=\alpha_Sb_{S,j},
\qquad
(P_\theta Q_\alpha)(\{S\})
=
\sum_{j=0}^{3}\pi_{\theta,j}\alpha_Sb_{S,j}.
\]
In particular,
\[
\alpha_S=0
\quad\Longrightarrow\quad
Q_\alpha(\{S\}\mid x_j)=0\ \text{for every }j
\quad\Longrightarrow\quad
(P_\theta Q_\alpha)(\{S\})=0.
\]
Thus
\[
\mathcal S_{\mathrm{out}}\subseteq\operatorname{supp}(\alpha),
\qquad
\kappa(Q_\alpha)
=
|\mathcal S_{\mathrm{out}}|
\le|\operatorname{supp}(\alpha)|
\le5.
\]
Together with feasibility and the profiled equality from the preceding
step, this proves the claim.
% lean: CausalSmith.Stat.LdpAteEfficiencySurface.outputCardinality_staircase_le_support_card
\end{enumerate}
\end{proof}
